如果一个集合~$A$ 是有限的，则它的大小可以用一个自然数来度量，而比较大小只需比较这些数即可。
若集合是无穷的，情况就复杂了。
无穷的第一个层级是\emph{可数无穷}集合。
如果集合~$A$ 的元素可以排列成一个\emph{枚举}——一个单向无穷的列表，即存在一个满射函数~$f\colon \PosInt \to A$，则称~$A$ 是\emph{可数的}。
如果一个集合可数但非有限，则它是\emph{可数无穷的}。
Cantor（康托尔）的\emph{之字形方法}表明，由可数无穷集合的元素对构成的集合也是可数的；
这可以用来证明，甚至有理数集~$\Rat$ 也是可数的。

然而，存在不可数的无穷集合：这样的集合称为\emph{不可数的}。
证明一个集合不可数有两种方法：直接使用\emph{对角论证}，或者通过\emph{归约}。
要给出对角论证，我们假设所讨论的集合~$A$ 是可数的，并利用一个假想的枚举来定义~$A$ 的一个元素；
按照我们的定义方式，该元素必定不同于枚举中的每一个元素。
因此，该枚举终究不能枚举~$A$ 的全部元素，从而我们证明了~$A$ 不可能存在枚举。
而归约则通过以满射的方式将~$A$ 的每个元素与某个已知的不可数集~$B$ 的元素关联起来，来证明~$A$ 是不可数的。
若能做到这一点，则~$A$ 的一个假想枚举就会产生~$B$ 的一个枚举。
由于~$B$ 是不可数的，~$A$ 不可能存在枚举。

一般而言，无穷集合的大小是可以比较的：如果~$A$ 和~$B$ 之间存在一个双射，则称它们大小相同，或\emph{等势}。
我们还可以定义，如果存在一个从~$A$ 到~$B$ 的单射函数，则~$A$ 不大于~$B$（记作~$\cardle{A}{B}$）。
根据 Schr\"oder-Bernstein 定理，这实际上确立了无穷集合之间按大小比较的序关系。
最后，\emph{康托尔定理}指出，对任意~$A$，都有~$\cardless{A}{\Pow{A}}$。
这推广了我们关于~$\Pow{\PosInt}$ 不可数的结果，并表明无穷的层级不止两个，而是有无穷多个。